\documentclass[10pt,a4paper]{amsart}
\usepackage[margin=19mm]{geometry}
\usepackage{amsmath,amssymb,mathtools}
\usepackage[scr=rsfs,cal=euler]{mathalfa}
\usepackage{xfrac}
\usepackage[unicode,hidelinks,bookmarks=false]{hyperref}
\newcommand{\ie}{i.e.\ }
\newcommand{\cf}{cf.\ }
\newcommand{\resp}{respectively\ }
\newcommand{\Subsection}{\subsection}
\providecommand{\nobreakdash}{\nobreak\hbox{-}}
\setlength{\emergencystretch}{2em}
\multlinegap0pt
\usepackage{mathtools}
\usepackage{amssymb}
\usepackage{algorithm}\usepackage{algpseudocode}
\newtheorem{lemma}{Lemma}
\newtheorem{theorem}{Theorem}
\title{DTPFI: A stable algorithm for recovering nonlinear energy potentials in phase field systems}
\author{Tianhao Ni, Jun Lai}
\date{Source: arXiv:2606.18607v1 (CC BY 4.0)}
\begin{document}
\maketitle

\noindent\emph{Source and licence.} DTPFI: A stable algorithm for recovering nonlinear energy potentials in phase field systems. Version arXiv:2606.18607v1 (CC BY 4.0), distributed by arXiv under the Creative Commons Attribution 4.0 International licence, http://creativecommons.org/licenses/by/4.0/, as recorded in the arXiv licence metadata for this exact version and shown on the abstract page https://arxiv.org/abs/2606.18607v1. Full licence notice: \texttt{licenses/arxiv-2606.18607-CC-BY-4.0.txt}.\par\smallskip

\noindent\textit{Editorial scope.} Counted unit: the article's theorem theorem (source0.tex lines 339--345). The article's complete original proof of it is delivered with it (source0.tex lines 347--374, one \texttt{proof} environment of 28 lines). No mathematics is written, repaired or re-derived: the statement and the proof are the article's own lines, in the article's own notation, and no retained line is substituted or re-flowed.  Retained uncounted as the source-local context the proof needs: lines 188--193; lines 209--227.  Nothing was omitted from any retained span, so the package carries no omission record.  No retained line contains a citation, so the wrapper carries no bibliography and no bibliography entry.  No retained line contains a figure environment, an image-inclusion command or drawing code, so no image is delivered and the package declares no asset dependency.  Editorial formatting only, in addition: an AMS \texttt{amsart} preamble that restates the article's own declarations and macros used by the retained lines, namely the environments \texttt{lemma}, \texttt{theorem} on separate counters, together with the standard packages those lines need.  The retained environments print in this document as Lemma 1, Theorem 1 in order of appearance, whereas the article prints the same items with the numbering of its own full document.  The complete native source of the article as distributed in the arXiv e-print for this exact version is bundled unchanged as \texttt{sources/203092/source0.tex}.
\begin{equation}
	\label{equ:opt2}
	\begin{aligned}
		\min\limits_{\mathbf{a} \in \mathbb{R}^{N+1}} \mathcal{J}_N(\mathbf{a}) = & \|u_{\mathbf{a}}(\cdot,t_2) - u_{\text{obs}}(\cdot,t_2)\|_{L^2(\Omega)}^2 +  \mathcal{R}(\mathbf{a}), 
	\end{aligned}
\end{equation}

\begin{lemma}
	\label{lemma:well_posedness_four_order}
	Consider the following initial-boundary value problem for $u(x,t)$:
	\begin{equation}
		\label{equ:linear CH}
		\left\{
		\begin{aligned}
			&u_{t} + a\Delta^2 u - \Delta (b(x,t)u) = p(x,t) && \text{in }  \Omega \times (t_1,t_2], \\
			&u(x,t_1) = \psi(x) && \text{in } \Omega, \\
			&u \text{ satisfies PBCs} && \text{on } \partial\Omega \times [t_1,t_2].
		\end{aligned}
		\right.
	\end{equation}
	Assume $a > 0$ and $b \in L^{\infty}(t_1,t_2; H_p^2(\Omega))$. For any source term $p \in L^{\infty}(t_1,t_2; L_p^{2}(\Omega))$ and initial value $\psi \in H_p^2(\Omega)$, there exists a unique solution $u \in L^{\infty}(t_1,t_2; H_p^2(\Omega)) \cap L^2(t_1,t_2; H^4_p(\Omega))$. Furthermore, the solution satisfies the stability estimate:
	\begin{align*}
		\|u\|_{L^{\infty}(t_1,t_2; H_p^2(\Omega))} + \|u\|_{L^2(t_1,t_2; H^4_p(\Omega))} \leq C \left( \|\psi\|_{H_p^2(\Omega)} + \|p\|_{L^{\infty}(t_1,t_2; L_p^{2}(\Omega))} \right),
	\end{align*}
	where the constant $C$ depends on $\Omega, T, a$, and $\|b\|_{L^{\infty}(t_1,t_2; H_p^2(\Omega))}$.
\end{lemma}

\begin{theorem}[Local strong convexity]
	Consider the optimization problem \eqref{equ:opt2} with the $\ell^2$ regularization term $\mathcal{R}(\mathbf{a}) = \beta\|\mathbf{a}\|_2^2$. Suppose the parameter estimate $\mathbf{a}$ lies within a $\rho_1$-neighborhood of the optimal solution $\mathbf{a}^{*}$, i.e., $\|\mathbf{a} - \mathbf{a}^{*}\|_2 \leq \rho_1$, and the corresponding residual satisfies:
	\begin{align}
		\|u_{\mathbf{a}}(\cdot,t_2) - u_{\text{obs}}(\cdot,t_2)\|_{L^2(\Omega)} \leq \rho_2.
	\end{align}
	Then, for a sufficiently large regularization weight $\beta > C\rho_2$, where $C$ is a constant determined by the stability of the second order sensitivity functions, the objective function $\mathcal{J}_N$ is strongly convex within this neighborhood.
\end{theorem}

\begin{proof}
	The Hessian matrix of the objective function $\mathcal{J}_N(\mathbf{a})$ is given by the sum of three terms:
	\begin{align*}
		H(\mathcal{J}_N) = \underbrace{2 \int_{\Omega} \mathbf{v}(\cdot,t_2) \mathbf{v}^T(\cdot,t_2) \mathrm{d}\mathbf{x}}_{G} + \underbrace{2 \int_{\Omega} (u_{\mathbf{a}}(\cdot,t_2) - u_{\text{obs}}(\cdot,t_2)) \mathbb{W}(\cdot,t_2) \mathrm{d}\mathbf{x}}_{E} + 2\beta I,
	\end{align*}
	where $\mathbf{v} = (v_0, \dots, v_N)^T$ is the vector of first-order sensitivities, $\mathbb{W} = (w_{ij})$ is the matrix of second order sensitivities, and $2\beta I$ is the Hessian of the regularization term $\beta \|\mathbf{a}\|_2^2$.
	
	We analyze the positive definiteness of $H(\mathcal{J}_N)$ by considering the quadratic form $\mathbf{h}^T H(\mathcal{J}_N) \mathbf{h}$ for an arbitrary non-zero vector $\mathbf{h} \in \mathbb{R}^{N+1}$:
	\begin{enumerate}
		\item The first term $G$ is the Gauss-Newton matrix, which is positive semi-definite: $\mathbf{h}^T G \mathbf{h} = 2 \|\mathbf{h}^T \mathbf{v}(\cdot,t_2)\|_{L^2(\Omega)}^2 \geq 0$.
		\item For the second term $E$, the quadratic form is 
		\begin{align*}
			\mathbf{h}^T E \mathbf{h} = 2 \int_{\Omega} (u_{\mathbf{a}}(\cdot,t_2) - u_{\text{obs}}(\cdot,t_2)) \left(\sum_{i,j=0}^N w_{ij} h_i h_j\right) \mathrm{d}\mathbf{x}.
		\end{align*}
		Let $\mathcal{W}_h = \sum_{i,j=0}^N w_{ij}(\cdot,t_2) h_i h_j$. Based on the stability estimate from Lemma \ref{lemma:well_posedness_four_order}, there exists a constant $C > 0$ such that $\|\mathcal{W}_h\|_{L^2(\Omega)} \leq C \|\mathbf{h}\|_2^2$. Applying the Cauchy-Schwarz inequality, we have:
		\begin{align*}
			|\mathbf{h}^T E \mathbf{h}| \leq 2 \|u_{\mathbf{a}} - u_{\text{obs}}\|_{L^2(\Omega)} \|\mathcal{W}_h\|_{L^2(\Omega)} \leq 2 \rho_2 C \|\mathbf{h}\|_2^2.
		\end{align*}
		This implies $E \succeq -2C\rho_2 I$.
		\item The regularization term contributes $2\beta I$, which is strictly positive definite for $\beta > 0$.
	\end{enumerate}
	
	Combining these bounds, we obtain:
	\begin{align*}
		H(\mathcal{J}_N) = G + E + 2\beta I \succeq 0 - 2C\rho_2 I + 2\beta I = 2(\beta - C\rho_2) I.
	\end{align*}
	By choosing $\beta > C\rho_2$, the eigenvalue of the Hessian is strictly positive, i.e., $H(\mathcal{J}_N) \succ 0$. Thus, $\mathcal{J}_N(\mathbf{a})$ is strongly convex in the specified $\rho_1$-neighborhood.
\end{proof}

\end{document}
